\chapter{Una computadora de células vivas}
% versión completa: chapters/21_living.tex

\pencilfig{21}{\pencilart{21}}{Neuronas vivas sobre un chip: debajo hay miles de electrodos, y cada uno sabe escuchar y hablar.}

¿Se puede construir una máquina no de silicio, sino de neuronas vivas? Se puede, y ya se construyen. Las neuronas se cultivan directamente sobre un chip con miles de electrodos, en una capa plana o en bolitas parecidas a un cerebro diminuto. Cada electrodo sabe escuchar y también enviar una señal. Es un banco de pruebas con el circuito abierto: todo se puede medir y en todo se puede intervenir. Solo que no tiene emisoras de radio: es un bus de datos sin registros de control.

\section*{Lo que la red hace sola}

Librado a sí mismo, el cultivo de vez en cuando estalla entero: casi todas las neuronas hacen clic juntas y luego se calman, con intervalos de un segundo a varios minutos. Es el oscilador ya conocido: la red se excita a sí misma y se cansa hasta que las conexiones reponen su reserva de sustancia.

\pencilfig{129}{%
% spike raster of 8 neurons in a culture left to itself: sparse single clicks and shared bursts
\foreach \b in {0.9,3.1,4.0,6.6}{\fill[pcAmber, opacity=0.18] (\b-0.1,-0.15) rectangle (\b+0.32,2.95);}
\foreach \y/\ss in {0/{0.96,1.0,1.13,2.43,3.0,3.12,3.24,4.02,4.09,4.15,6.66,6.69,6.81},
  1/{0.46,0.88,1.0,1.05,3.09,3.2,3.94,4.04,4.37,6.65,6.71},
  2/{0.73,0.84,0.94,1.41,3.08,3.12,3.21,4.05,4.06,4.17,6.59,6.63},
  3/{0.89,0.93,1.06,3.1,3.19,3.28,3.89,3.94,4.13,4.2,4.31,6.63,6.65,6.78},
  4/{0.87,0.96,3.09,3.15,3.23,3.73,3.96,4.06,4.19,5.98,6.66,6.68,6.75},
  5/{0.44,0.92,0.97,3.05,3.22,3.27,4.01,4.07,4.17,5.76,6.57,6.73,6.75},
  6/{0.92,1.0,1.73,2.85,3.18,3.19,4.0,4.07,6.53,6.66,6.73},
  7/{0.87,0.92,3.13,3.13,3.27,3.99,4.06,4.17,6.44,6.61,6.72,7.13}}{
  \foreach \s in \ss {\draw[pcBlue, line width=0.7pt] (\s,0.35*\y) -- (\s,0.35*\y+0.25);}}
\draw[arr] (0,-0.3) -- (7.5,-0.3); \node[lbl, anchor=north] at (3.75,-0.32) {tiempo};
\node[lbl, anchor=east, align=right] at (-0.05,1.35) {ocho\\neuronas};
\node[lbl, text=pcAmber!70!black, anchor=south] at (3.55,2.95) {estallidos conjuntos};
}{Un cultivo sobre un chip, librado a sí mismo. Los clics sueltos son raros, y de vez en cuando casi toda la red estalla a la vez y se calma hasta que las conexiones reponen su reserva.}

\section*{Un maestro sin dopamina}

¿Cómo se le enseña a una red así si no hay dopamina? Los primeros métodos fueron eléctricos. Por ejemplo: estimular la red hasta que dé la respuesta deseada, y parar en seguida. La respuesta se afianza. En otro experimento, el cultivo controlaba un sencillo “cuerpo” virtual y aprendía a llevarlo hacia una meta. En todos estos experimentos, el maestro es un patrón de corriente que elige un ingeniero.

\section*{Dopamina con un destello de luz}

Una vez se intentó que el maestro fuera químico. En una de las plataformas con bolitas de neuronas vivas se añadió al líquido nutritivo dopamina dentro de una “jaula” fotosensible que le impide actuar. Cuando la red respondía a un estímulo con más fuerza que antes, se la “premiaba”: un destello ultravioleta rompía la jaula y la dopamina quedaba libre. A lo largo de doscientas cuarenta repeticiones, las respuestas en conjunto crecieron. Pero los propios autores escriben con honestidad que con un solo experimento así no se puede afirmar que la causa sea la dopamina.

\section*{Una pértiga sobre un carrito}

En un experimento reciente se enseñó a bolitas de neuronas a equilibrar una pértiga sobre un carrito, un problema clásico para los programas de aprendizaje. Qué pulsos de entrenamiento aplicar lo elegía un programa. Las señales elegidas por el programa funcionaban mejor que las aleatorias, pero tras tres cuartos de hora de descanso lo aprendido desaparecía. Y si se bloqueaban las conexiones \nt{GLUT}, no había aprendizaje en absoluto: lo aprendido vive justo ahí. El maestro, en cambio, sigue estando afuera; solo que ahora no es un ingeniero, sino un programa.

\section*{Un reservorio vivo}

Hay otro enfoque: a la red viva no se le enseña nada. Se la usa como “reservorio”, un sistema rico y enmarañado cuyas respuestas aprende a descifrar un simple contador electrónico. Así, una bolita viva ayudó a distinguir voces. Eso sí, según los autores, las conexiones de la propia bolita también cambiaban durante el trabajo.

\section*{El precio y la pregunta}

Un millón de neuronas sobre un chip consume más o menos una diezmilésima de vatio, mucho menos que la electrónica que las rodea. Y el banco de pruebas no sabe decidir por sí solo qué cuenta como éxito: ese papel siempre lo cumple alguien de afuera. Y a medida que estos sistemas crecen, surge también una cuestión ética que nadie ha resuelto todavía.

\fullversion{21}
